\documentclass[quick,pro]{debate}
\motion{对宠物的爱，爱的更是它 / 爱的更是我}
\stance{正方}{爱的更是它}
\begin{document}
\debatetitle
\begin{thesisbox}[一句话立论]养宠这件事里真正需要解释的，不是我们得到了什么，而是我们在得不到任何东西的时候还在为它做什么；答案落在它身上，所以爱的更是它。\end{thesisbox}
\section{定义与判准（标准表述）}
\begin{fields}
\field{爱}{对某个对象有深挚的感情，以及由这份感情带来的持续关切与行动。（有感受，也有为它做的事）
}
\field{更}{这是一道比较题，两边都在，比的是分量。我方承认养宠让人快乐、让人被需要。
}
\field{它}{你家那一只，有自己的脾气和需要，不随你的想法改变。
}
\field{判准（中立，两条通道）}{一、这份爱的内容是被谁塑造的；二、两边的需要真正撞上时，这份爱跟着谁走。
\begin{bullets}
\item 我方要证明：这份感情主要由它塑造，冲突时多数人让的是自己。
\item 对方要证明：反过来。
\end{bullets}
}
\field{对方提偏向定义或判准时的 30 秒口径}{如果对方把爱定义成“满足自我需要的情感投入”，立刻指出这让母爱、爱情、友情全都“更是我”，命题不可证伪；然后问那一问：“按您的标准，有没有任何一种爱是更指向对方的？”如果对方只肯要“冲突时刻”这一条通道，回答：冲突时刻是必要检验，不是唯一检验。
}
\end{fields}
\section{我方论点}
\begin{longtable}{@{}L{\dimexpr0.044\linewidth-1.50\tabcolsep}L{\dimexpr0.165\linewidth-1.50\tabcolsep}L{\dimexpr0.438\linewidth-1.50\tabcolsep}L{\dimexpr0.353\linewidth-1.50\tabcolsep}@{}}
\toprule \thead{标签} & \thead{一句话} & \thead{核心数据 / 学理 / 案例} & \thead{出处}\\\midrule\endhead
让路的是我 & 冲突真的发生时，多数人让的是自己，而且让得最彻底的地方恰恰是自己拿不到好处的地方 & 学理：照料行为系统的目标是对方的福祉，与依恋系统不是一回事。数据：2025 年城镇犬猫消费 3126 亿元，医疗占 27.6\%。案例：2020 年武汉封城，8 名核心工作人员与约 60 名志愿者救助约 5000 只留守宠物 & Mikulincer \& Shaver 2005；《2026 年中国宠物行业白皮书》（央视网 2026-01-05）；谷雨工作室 2020-02-28\\
认的是这一只 & 这份感情围着它的具体脾气长出来，换一只同款就不成立 & 学理：历史关系属性（Whiting 1991、Delaney 1996）。数据：29 对犬与主人的陌生情境实验，狗找的是主人不是陌生人，狗与主人皮质醇正相关。案例：71.6\% 的年轻饲主困扰“它老了怎么办” & SEP「Love」；Ryan 等 2019, Front Behav Neurosci 13:162；青年志《2025 青年养宠观报告》\\
\bottomrule\end{longtable}
\section{对方论点 → 一句话反驳}
\begin{longtable}{@{}L{\dimexpr0.249\linewidth-1.33\tabcolsep}L{\dimexpr0.384\linewidth-1.33\tabcolsep}L{\dimexpr0.366\linewidth-1.33\tabcolsep}@{}}
\toprule \thead{对方会说} & \thead{我方一句话回} & \thead{对方追问时}\\\midrule\endhead
养宠起点是孤独、是陪伴需求 & 需要它和爱它是两件事，交朋友也是因为需要朋友 & 孤独推人去建立关系，不说明关系是假的\\
拟人化：你爱的是投射（Epley） & 投射不会被纠正，这份感情天天在被它纠正 & 我按它的反应改，不按我的想象坚持\\
短头品种是按我们的审美选的 & 那是买之前的审美，爱是进门之后的事 & 买的人和爱的人在时间上是两个阶段\\
喂胖了、穿衣服伤了它 & 无知推不出自私，做错事不等于没爱 & 父母逼你多穿一件、多吃一碗，效果也常常是错的\\
延迟安乐死：主人怕自己难过 & 拖的是我的不舍，签字的是它的痛苦 & 没有人从拖延里得到享受，得到的是折磨\\
弃养 & 用不爱它的人来定义爱，方法本身有问题 & 请给家庭弃养的占比，给不出就只是印象\\
朋友圈、写真是自我形象经营 & 晒的是这段关系，人也给孩子拍照 & 买玩具的接近九成，高于买服装配饰的一半多\\
Archer：宠物是社会寄生者 & 那是解释来源，不是判定指向，中间隔着发生学谬误 & 母爱也能被解释成基因策略\\
自我延伸：起作用的是“我变大了” & 自我延伸理论本来解释爱情，照这用法所有爱情都是自恋 & 它真的会拒绝：会咬你、会躲、会跳走\\
宠物是可控的关系 & 这是全世界最不可控的关系之一 & 它会用十五年的寿命离开你，你第一天就知道\\
\bottomrule\end{longtable}
\section{必问与必答}
\begin{fields}
\field{必问 （对方不答就点名）}{
\begin{enumerate}[leftmargin=1.8em,itemsep=2pt,topsep=2pt]
\item 按您方的标准，有没有任何一种爱是更指向对方的？（一辩稿抛出的那个问题）
\item 家庭弃养占城镇犬猫的多少？
\item 挑品种发生在认识它之前还是之后？
\item 你方那些依恋研究，因变量长在谁身上？
\item 临终止痛这件事，受益人是谁？
\end{enumerate}
}
\field{必答 （15 秒以内正面回答）}{
\begin{bullets}
\item “你承认养宠是为了陪伴吗？” → 承认。得到是副产品，指向还在它身上。
\item “84\% 的主人说自家狗不胖。” → 承认。那是认知偏差，不是动机，人对自己的体重也估不准。
\item “安乐死为什么大家都拖？” → 拖的是我的不舍，签字的是它的痛苦。
\item “那件小裙子是给谁穿的？” → 我方也反对那种做法。一年一次的写真和一天两顿的饭，哪个更能说明日常。
\item “你方承不承认自己也获益？” → 承认，第一句就说了。
\end{bullets}
}
\end{fields}
\section{自由辩战场概要}
\begin{longtable}{@{}L{\dimexpr0.123\linewidth-1.50\tabcolsep}L{\dimexpr0.318\linewidth-1.50\tabcolsep}L{\dimexpr0.256\linewidth-1.50\tabcolsep}L{\dimexpr0.303\linewidth-1.50\tabcolsep}@{}}
\toprule \thead{战场} & \thead{开场问题} & \thead{目标} & \thead{收束语}\\\midrule\endhead
有和更之间的数字 & 家庭弃养占城镇犬猫的比例是多少？ & 让评委看到对方拿不出“更”的比例 & 反方举的是“有人做过”，我方问的是“多数人往哪边”\\
投射会不会被纠正 & 你以为猫想被抱，它挠了你，你会不会改？ & 把拟人化从“对象是假的”拉回“理解的方式” & 真正的自我投射不会被对方修正\\
临终那一段 & 一只肾衰竭的老猫做皮下补液，这笔钱买到了什么？ & 给评委一个画面收束全场 & 我方承认得到，我方只说得到是副产品\\
\bottomrule\end{longtable}
时间分配：前 90 秒战场一，抛出必问一到三；中 90 秒战场二，处理对方进攻；后 60 秒战场三，不开新战场，留 20 秒备用。站位：一辩守定义与判准，二辩接驳论过的点，三辩接盘问成果，四辩收束与价值。

\section{如果……就……}
\begin{contingency}
\ifthen{对方定义不同}{用“命题不可证伪”这一击，再把中立判准两条通道复述一遍。}
\ifthen{对方拿出没预料到的数据}{先问机构、年份、样本，再问它测的因变量在谁身上。}
\ifthen{论点二被打穿（对方拿 Badhwar 打不可替代）}{收缩到论点一，用“冲突时让路的是我”扛全场。}
\ifthen{对方回避必问}{点名、重复、不换题。第二次不答就说“这个问题对方已经两次没有回答”。}
\end{contingency}
\section{全场三条铁律（背下来）}
\begin{enumerate}[leftmargin=1.8em,itemsep=2pt,topsep=2pt]
\item 不能说“我从养宠里什么都没得到”。
\item 不能说“所以它也爱我”，要说“它是一个会区分人的真实他者”。
\item 不能用“我为它花了多少钱”当爱它的证据，一说对方就问“买裙子算不算”。
\end{enumerate}
\section{各环节发言人与时长}
\begin{longtable}{@{}L{\dimexpr0.201\linewidth-1.50\tabcolsep}L{\dimexpr0.201\linewidth-1.50\tabcolsep}L{\dimexpr0.201\linewidth-1.50\tabcolsep}L{\dimexpr0.396\linewidth-1.50\tabcolsep}@{}}
\toprule \thead{环节} & \thead{我方发言人} & \thead{时长} & \thead{一句话提醒}\\\midrule\endhead
1 开篇立论 & 正一 & 180 秒 & 开头 20 秒把定义与判准讲完\\
2 被质询 & 正一 & 120 秒双边 & 承认常识，守住“得到是副产品”\\
4 质询反一 & 正四 & 120 秒双边 & 先打定义杀边界那一链\\
7 驳论 & 正二 & 120 秒 & 留 100 字临场位回应反二\\
8 对辩 & 正二 & 90 秒 & 回应加反问，说完就停\\
10 盘问 & 正三 & 90 秒 & 同一问题先问二辩再问四辩\\
11 被盘问 & 正二、正四 & 90 秒 & 两人答案必须一致\\
12 盘问小结 & 正三 & 90 秒 & 预判反三会怎么总结，提前压住\\
15 自由辩 & 全员 & 240 秒 & 三个战场按时间推进\\
18 结辩 & 正四 & 210 秒 & 留 30–40 秒真回应反四，不背稿\\
\bottomrule\end{longtable}
\end{document}
